\begin{trapbox}

	\begin{description}[style=nextline, leftmargin=2.8em, labelsep=0.5em, itemsep=0.35em]
		\item[\textbf{\textcolor{CTrap}{易错点一（忽略重复计数）}}]
			见到 $\{1,2,2,3\}$ 误以为元素个数是 $4$，直接按 $n=4$ 计算子集数。
		\item[\textbf{\textcolor{CTrap}{正确理解（互异化简）：}}]
			集合元素必须互异，重复书写的元素只能算一个。应先将集合化简为 $\{1,2,3\}$，再认定 $n=3$。
		\item[\textbf{\textcolor{CTrap}{错因分析（惯性思维）：}}]
			受到日常书写习惯影响，把列举法中的“逗号分隔项数”直接当作元素个数，忘记集合的互异性约束。
	\end{description}

	\medskip
	\begin{description}[style=nextline, leftmargin=2.8em, labelsep=0.5em, itemsep=0.35em]
		\item[\textbf{\textcolor{CTrap}{易错点二（参数不检验互异）}}]
			在条件 $1\in A$ 下解出 $a$ 的若干候选值后，直接全部保留作为最终答案。
		\item[\textbf{\textcolor{CTrap}{正确理解（参数代回检验）：}}]
			必须将每个候选值代回原集合，检查是否出现相同元素。若造成元素重复，该参数值应舍去。
		\item[\textbf{\textcolor{CTrap}{错因分析（目标单一）：}}]
			只关注“$1$ 属于集合”这一条件，忽略了集合定义本身自带的互异性要求，没有养成多条件综合检验的习惯。
	\end{description}

	\medskip
	\begin{description}[style=nextline, leftmargin=2.8em, labelsep=0.5em, itemsep=0.35em]
		\item[\textbf{\textcolor{CTrap}{易错点三（混淆顺序意义）}}]
			认为 $\{a,b\}$ 与 $\{b,a\}$ 是两个不同的集合，尤其在用参数表示元素时认为顺序蕴含某种对应关系。
		\item[\textbf{\textcolor{CTrap}{正确理解（无序性等价）：}}]
			二集合相等只看所含元素是否完全相同，与书写次序无关。因此 $\{a,b\}=\{b,a\}$ 恒成立，无需考虑顺序匹配。
		\item[\textbf{\textcolor{CTrap}{错因分析（序列观念干扰）：}}]
			将高中数学中的“集合”与日常的“序列”、“坐标”等概念混淆，把元素的排列顺序误读为集合的一部分属性。
	\end{description}

\end{trapbox}
